Ein Körper bewegt sich im Gegenuhrzeigersinn gleichförmig auf einer Kreisbahn. Zu Beginn beträgt seine Geschwindigkeit \vkO\ im tiefsten Punkt der Bahn. Nach \tkO\ hat der horizontale Teil der Geschwindigkeit auf \vxeO\ abgenommen.
\begin{abcliste}
\abc Konstruieren Sie zeichnerisch den durchlaufenen Drehwinkel.
\abc Berechnen Sie den Radius der Kreisbahn.
\end{abcliste}
